\documentclass[11pt,a4paper]{article}

\usepackage[margin=2.5cm]{geometry}
\usepackage[T1]{fontenc}
\usepackage[utf8]{inputenc}
\usepackage{mathptmx}
\usepackage{microtype}
\usepackage{amsmath,amssymb}
\usepackage{booktabs}
\usepackage{enumitem}
\usepackage{graphicx}
\usepackage{caption}
\usepackage{float}
\usepackage[hidelinks]{hyperref}

\setlist{nosep}
\captionsetup{font=small,labelfont=bf}

% ---------------------------------------------------------------
% small macros used in every exercise
% ---------------------------------------------------------------
\newcommand{\vx}{\mathbf{x}}
\newcommand{\vs}{\mathbf{s}}
\newcommand{\ve}{\mathbf{e}}
\newcommand{\A}{\mathbf{A}}
\newcommand{\M}{\mathbf{M}}
\newcommand{\Smat}{\mathbf{S}}
\newcommand{\Vone}[1]{V_1(#1)}
\newcommand{\sumcheck}[1]{\begin{center}\small (sum $= #1$)\end{center}}

% Exercise heading: number, title, author
\newcommand{\exheading}[3]{%
  \clearpage
  \section*{Exercise #1: #2}
  \addcontentsline{toc}{section}{Exercise #1: #2}
  %\begin{flushright}\small\textit{Solved by #3}\end{flushright}
  \begin{flushright}\small\textit{}\end{flushright}
  
  }
\newcommand{\exblock}[1]{\subsection*{#1}}

\title{\textbf{The PageRank Algorithm}\\[0.4em]
  \large Homework 1 and Homework 3\\[0.3em]
  \normalsize Computational Linear Algebra for Large Scale Problems\\
  Politecnico di Torino, A.Y.~2025/2026}
\author{Ömer Özkan (s360761) \and Zeynep Sude Ulaş (s364189)}
\date{September 2026}

\begin{document}

\maketitle

\begin{abstract}
\noindent
This report is our work on the paper \emph{The \$25,000,000,000 Eigenvector: The Linear Algebra Behind Google} by K.~Bryan and T.~Leise.
Part~I is Homework~1: we describe the PageRank algorithm, we explain why it works and when it does not, we present our own implementation (a hand-written sparse matrix and the power method), we test it on the two small webs of the paper and on the Hollins dataset (6012 pages), and we solve Exercises~11 and~17.
Part~II is Homework~3: we solve ten more exercises of the paper (1, 2, 3, 5, 6, 7, 8, 9, 12 and 13).
All numbers in this report are produced by the code in our repository.
\end{abstract}

\tableofcontents

\clearpage
\part{Homework 1: The PageRank Project}
\section{Introduction}

\subsection*{The task}

The homework asks us to read the paper of Bryan and Leise~\cite{bryanleise}, to write our own code for the PageRank algorithm, to test it on the two webs of the paper (Figure~2.1 and Figure~2.2), to test it on the given dataset, and to solve two exercises of our choice.
We chose \textbf{Exercise~11} (a new page is added to the web and the ranking is computed with the matrix $\M$) and \textbf{Exercise~17} (how the damping factor $m$ should be chosen).

\subsection*{The dataset}

The dataset is the \emph{Hollins} web graph: 6012 pages of the website of Hollins University and 23\,875 links between them.
It is stored in the file \texttt{hollins.dat}.

\subsection*{Notation}

We use the notation of the paper.
A web has $n$ pages, numbered $1,\dots,n$.
The vector $\vx = (x_1,\dots,x_n)^T$ contains the importance scores, and we always scale it so that $\sum_i x_i = 1$.
$\Vone{\A}$ is the eigenspace of the matrix $\A$ for the eigenvalue~$1$.
When we quote a formula of the paper we use the paper's number, for example ``formula (2.1)'' or ``formula (3.2)''.

\clearpage
\section{The PageRank Algorithm}

\subsection{How it works}

\subsubsection*{The idea}

A link from page $j$ to page $k$ is a \emph{vote} of page $j$ for page $k$.
A vote from an important page should count more than a vote from an unimportant page.
Also, a page should not become powerful only by linking to many pages: each page has one vote in total, and it divides this vote equally among its outgoing links.

Let $L_k$ be the set of pages that link to page $k$ (the \emph{backlinks} of $k$), and let $n_j$ be the number of outgoing links of page $j$.
The paper defines the score of page $k$ by formula (2.1):
\[
  x_k = \sum_{j \in L_k} \frac{x_j}{n_j}.
\]

\subsubsection*{The link matrix $\A$}

The $n$ equations of formula (2.1) can be written as one matrix equation $\A\vx = \vx$, where
\[
  A_{ij} = \begin{cases} 1/n_j & \text{if page $j$ links to page $i$,}\\ 0 & \text{otherwise.}\end{cases}
\]
So column $j$ of $\A$ describes the outgoing links of page $j$, and row $i$ describes the backlinks of page $i$.
Every column sums to one (page $j$ divides its whole vote), so $\A$ is \emph{column-stochastic}.
The score vector $\vx$ is an eigenvector of $\A$ with eigenvalue~$1$.

\begin{itemize}
  \item For the four-page web of Figure~2.1 the paper finds $\vx = \frac{1}{31}(12,\,4,\,9,\,6)^T$, so page~1 is the most important.
  \item Every column-stochastic matrix has $1$ as an eigenvalue (Proposition~1 of the paper), so a solution always exists when there are no dangling nodes.
\end{itemize}

\subsubsection*{The matrix $\M$}

The matrix $\A$ has two problems (see Section~\ref{sec:when}).
To fix them the paper mixes $\A$ with the matrix $\Smat$ whose entries are all $1/n$, using a number $m \in [0,1]$ called the \emph{damping factor}:
\[
  \M = (1-m)\A + m\Smat \qquad\text{(formula (3.1))}.
\]
Google used $m = 0.15$, and we use the same value.
Because $\Smat\vx = \vs$ when $\sum_i x_i = 1$, where $\vs$ is the vector with all entries $1/n$, the equation $\M\vx = \vx$ can also be written as
\[
  \vx = (1-m)\A\vx + m\vs \qquad\text{(formula (3.2))}.
\]
Formula (3.2) is the form we use in the code: we never build $\M$, we only multiply by the sparse matrix $\A$ and add the constant vector $m\vs$.

\subsubsection*{The power method}

For a web with billions of pages we cannot compute eigenvectors directly.
The paper uses the \emph{power method}: start from any positive vector $\vx_0$ with $\|\vx_0\|_1 = 1$ and repeat
\[
  \vx_{k+1} = \M\vx_k = (1-m)\A\vx_k + m\vs .
\]
The sequence $\vx_k$ converges to the score vector $\mathbf{q}$ (Theorem~4.2 of the paper).
We stop when two consecutive vectors are close in the 1-norm, $\|\vx_{k+1}-\vx_k\|_1 < \texttt{tol}$.
Each step costs one product $\A\vx$, which is cheap because $\A$ is sparse: on Hollins only $0.066\%$ of the entries of $\A$ are non-zero.

\subsection{Why it works}

The proofs are in Section~3.2 and Section~4 of the paper. We summarise the chain of ideas.

\begin{enumerate}
  \item \textbf{$\M$ is column-stochastic} (Exercise~7). So $1$ is an eigenvalue of $\M$ and $\Vone{\M}$ is not empty.
  \item \textbf{$\M$ is positive.} Every entry of $\Smat$ is $1/n > 0$, so for $m > 0$ every entry $M_{ij} \geq m/n > 0$.
  \item \textbf{A positive column-stochastic matrix has a one-dimensional $V_1$} (Lemma~3.2).
  Any eigenvector in $\Vone{\M}$ has all positive or all negative entries (Proposition~2), and two independent vectors could be combined into a vector with mixed signs (Proposition~3), which is impossible.
  So there is exactly one score vector $\mathbf{q}$ with positive entries and $\sum_i q_i = 1$.
  \item \textbf{The power method converges to $\mathbf{q}$} (Proposition~5).
  Write $\vx_0 = \mathbf{q} + \mathbf{v}$ with $\sum_j v_j = 0$.
  Then $\M^k\vx_0 = \mathbf{q} + \M^k\mathbf{v}$, and Proposition~4 gives $\|\M^k\mathbf{v}\|_1 \le c^k\|\mathbf{v}\|_1$ with $c < 1$.
  So the error goes to zero.
  \item \textbf{Speed.} The error decreases roughly like $|\lambda_2|^k$, where $\lambda_2$ is the second largest eigenvalue of $\M$, and the paper states that $|\lambda_2| \le 1-m$.
  So a larger $m$ gives faster convergence. We check this in Exercise~17.
\end{enumerate}

\subsection{When it does not work}
\label{sec:when}

\subsubsection*{Disconnected webs: the ranking of $\A$ is not unique}

If the web (seen as an undirected graph) has $r$ disconnected pieces, then $\dim\Vone{\A} \ge r$ (Section~2.2.1 of the paper).
Figure~2.2 has two pieces and $\dim\Vone{\A} = 2$: every combination of the two basis vectors is an eigenvector, so $\A$ does not tell us which one is ``the'' ranking.
The matrix $\M$ solves this because it is positive, so $\dim\Vone{\M} = 1$ for every web (Exercises~2, 3 and 13).

\subsubsection*{Dangling nodes: $\A$ loses mass}

A \emph{dangling node} is a page with no outgoing links.
Its column in $\A$ is all zeros, so $\A$ is only \emph{sub}-stochastic and $1$ may not be an eigenvalue at all.
The paper does not treat this case, but the Hollins dataset has 3189 dangling pages (53\% of all pages), so our code must handle it.
We use the standard fix: a dangling page gives its vote equally to all $n$ pages.
In the code this means adding the number $\frac{1}{n}\sum_{j \text{ dangling}} x_j$ to every entry of $\A\vx$.
After this fix the matrix is column-stochastic again and the theory above applies.

\subsubsection*{The value $m = 0$: the power method may not converge}

With $m = 0$ we have $\M = \A$, which is not positive.
Then the power method can fail.
On Figure~2.2 we ran the iteration $\vx_{k+1} = \A\vx_k$ from the starting vector $\vx_0 = (0.5,\,0.1,\,0.2,\,0.1,\,0.1)^T$: after 1000 steps the residual $\|\vx_{k+1}-\vx_k\|_1$ was still $1.0$, because $\A$ has the eigenvalue $-1$ and the iterates jump between two vectors.
With $m = 0.15$ the same start converges in 143 steps to $(0.2,\,0.2,\,0.285,\,0.285,\,0.03)^T$, the vector given in the paper.

\subsubsection*{The value $m$ close to $1$: the links are ignored}

When $m \to 1$, $\M \to \Smat$ and all scores tend to $1/n$.
The ranking then says nothing about the web.
The choice of $m$ is therefore a compromise between speed and information (Exercise~17).

\subsubsection*{Link spam}

Because a page can raise its score by creating pages that link to it, PageRank can be manipulated.
Exercise~1 shows a small example of this.

\clearpage
\section{Implementation and Results}
\label{sec:implementation}

The code is written in Python with NumPy only (no SciPy, no sparse library).
It is in our repository~\cite{repo}. The files are listed in Table~\ref{tab:files}.

\begin{table}[H]\centering
\begin{tabular}{ll}
\toprule
\textbf{File} & \textbf{Content} \\
\midrule
\texttt{toolkit.py} & the matrices, the sparse storage, the power method, the file reader \\
\texttt{webs.py} & every small web used in this report, written as a dictionary \\
\texttt{run\_paper.py} & tests against the numbers printed in the paper \\
\texttt{hollins\_main.py} & PageRank on the Hollins dataset \\
\texttt{ex01.py}, \dots, \texttt{ex17.py} & one script per computational exercise \\
\bottomrule
\end{tabular}
\caption{Files of the project}
\label{tab:files}
\end{table}

\subsection{Data structures}

\subsubsection*{A web is a dictionary}

A web is stored as a Python dictionary \verb|{page: [pages it links to]}| with pages numbered $1,\dots,n$ as in the paper.
For example Figure~2.1 is
\begin{verbatim}
FIG_2_1 = {1: [2, 3, 4], 2: [3, 4], 3: [1], 4: [1, 3]}
\end{verbatim}
A page with an empty list is a dangling node.

\subsubsection*{Dense matrix for small webs}

The function \texttt{build\_A} fills an $n \times n$ NumPy array with $A_{ij} = 1/n_j$.
We use it only for the small webs, to print $\A$ and $\M$ and to compute all eigenvalues with \texttt{numpy.linalg.eig} for cross-checking (for example to find $\dim\Vone{\A}$).
The dense matrix is \emph{not} used by the solver.

\subsubsection*{Sparse matrix (CSR) for the dataset}

For Hollins a dense $\A$ would need $6012^2 \cdot 8$ bytes $\approx 289$~MB, but only $23\,875$ entries are non-zero.
We store $\A$ in the \emph{compressed sparse row} format with three arrays, built by hand in \texttt{build\_csr}:
\begin{itemize}
  \item \texttt{AA}: the non-zero values of $\A$, read row by row;
  \item \texttt{JA}: the column index of each value in \texttt{AA};
  \item \texttt{IA}: for each row $i$, the position in \texttt{AA} where row $i$ starts (length $n+1$).
\end{itemize}
The non-zeros of row $i$ are \texttt{AA[IA[i]:IA[i+1]]}.
The three arrays for Hollins take $0.43$~MB, about 670 times less than the dense matrix.
The function \texttt{csr\_matvec} computes $\mathbf{y} = \A\vx$ row by row:
\[
  y_i = \sum_{k=\texttt{IA}[i]}^{\texttt{IA}[i+1]-1} \texttt{AA}[k]\; x_{\texttt{JA}[k]} .
\]
We tested that this product equals the dense product \texttt{A @ v} on a random vector (difference $0$).

\subsection{The solver}

The only solver of the project is the function \texttt{pagerank(links, n, m, tol)}.
It starts from $\vx_0 = (1/n,\dots,1/n)^T$ and repeats three steps until $\|\vx_{k+1}-\vx_k\|_1 < \texttt{tol}$:
\begin{enumerate}[label=\textbf{Step \Alph*.},leftmargin=*]
  \item $\mathbf{y} = \A\vx_k$ with the CSR product.
  \item $\mathbf{y} \leftarrow \mathbf{y} + \frac{1}{n}\sum_{j\ \text{dangling}} (\vx_k)_j$: the mass of the dangling pages is given back, equally to all pages.
  \item $\vx_{k+1} = (1-m)\mathbf{y} + m\vs$: this is formula (3.2).
\end{enumerate}
After Step~C the sum of $\vx_{k+1}$ is already $1$ in exact arithmetic; we divide by the sum once more only to remove rounding errors.
The function returns the vector $\vx$, the number of iterations $k$ and the last residual $\|\vx_{k+1}-\vx_k\|_1 = \|\M\vx_k - \vx_k\|_1$.
If the residual is not below \texttt{tol}, the method did not converge.

With $m = 0$ the same function iterates $\vx_{k+1} = \A\vx_k$, so we also use it to rank a web with formula (2.1) (Exercises~1 and~12).

\subsubsection*{Reading the dataset}

The function \texttt{load\_dat} reads \texttt{hollins.dat}: the first line gives $n$ and the number of links, then $n$ lines give the URL of each page, then one link per line as ``source target''.
Following the paper, a link from a page to itself would be ignored, and a repeated link would be counted once.
In the Hollins file there are no self-links and no repeated links.

\subsection{Tests on the two webs of the paper}

Before using the dataset we checked every number that the paper prints (script \texttt{run\_paper.py}).
Table~\ref{tab:paper-checks} shows that all of them agree.
For each value the allowed difference is one unit of the last printed decimal of the paper.

\begin{table}[H]\centering
\begin{tabular}{llcc}
\toprule
\textbf{Web} & \textbf{Quantity checked} & \textbf{Paper} & \textbf{Max.\ difference} \\
\midrule
Fig.~2.1 & eigenvector of $\A$ $= (12,4,9,6)/31$ & p.~3 & $1.7\cdot10^{-16}$ \\
Fig.~2.1 & $\dim\Vone{\A} = 1$ & p.~3 & exact \\
Fig.~2.1 & matrix $\M$ (four decimals) & Example 1 & $3.3\cdot10^{-6}$ \\
Fig.~2.1 & scores $0.368,\,0.142,\,0.288,\,0.202$ & Example 1 & $1.9\cdot10^{-4}$ \\
Fig.~2.2 & $\dim\Vone{\A} = 2$ & p.~4 & exact \\
Fig.~2.2 & matrix $\M$ of formula (3.3) & p.~6 & $5.6\cdot10^{-17}$ \\
Fig.~2.2 & scores $0.2,\,0.2,\,0.285,\,0.285,\,0.03$ & Example 2 & $5.6\cdot10^{-17}$ \\
\bottomrule
\end{tabular}
\caption{Checks of our code against the values printed in the paper}
\label{tab:paper-checks}
\end{table}


The ranking for Figure~2.1 with $\M$ is
\[
  \vx \approx [0.368151,\ 0.141809,\ 0.287962,\ 0.202078]
\]
\sumcheck{1.000}
and for Figure~2.2 it is exactly $[0.2,\ 0.2,\ 0.285,\ 0.285,\ 0.03]$, as in Example~2 of the paper.
The order of Figure~2.1 is $1 > 3 > 4 > 2$, the same order as with $\A$.

We also tested the dangling-node fix on a small web (Figure~2.1 without the link $3 \to 1$, so page~3 is dangling): the solver and a dense eigenvector computation on the corrected matrix give the same vector (difference $4\cdot10^{-14}$).

\subsection{Results on the Hollins dataset}
\label{sec:hollins}

Table~\ref{tab:hollins} summarises the run of \texttt{hollins\_main.py} with $m = 0.15$ and $\texttt{tol} = 10^{-12}$.

\begin{table}[H]\centering
\begin{tabular}{lr}
\toprule
\textbf{Quantity} & \textbf{Value} \\
\midrule
pages $n$ & 6012 \\
links (non-zeros of $\A$) & 23\,875 \\
fill of $\A$ & 0.066\% \\
dangling pages & 3189 (53\%) \\
memory, CSR arrays & 0.43 MB \\
memory, dense matrix & 289 MB \\
iterations ($m=0.15$, $\texttt{tol}=10^{-12}$) & 138 \\
final residual $\|\M\vx_k-\vx_k\|_1$ & $9.7\cdot10^{-13}$ \\
time (build CSR + iterations) & 1.0 s \\
\bottomrule
\end{tabular}
\caption{PageRank on the Hollins dataset}
\label{tab:hollins}
\end{table}


The most important page is the home page of the university, with a score about 120 times larger than the uniform score $1/n = 0.000166$.
Table~\ref{tab:hollins-top} lists the ten best pages.
The pages of the admissions area are high because many pages of the site link to them.

\begin{table}[H]\centering
\small
\begin{tabular}{crll}
\toprule
\textbf{Rank} & \textbf{Page} & \textbf{Score} & \textbf{URL} \\
\midrule
1 & 2 & 0.017669 & \texttt{hollins.edu/} \\
2 & 37 & 0.010238 & \texttt{hollins.edu/admissions/visit/visit.htm} \\
3 & 38 & 0.009457 & \texttt{hollins.edu/about/about\_tour.htm} \\
4 & 61 & 0.009069 & \texttt{hollins.edu/htdig/index.html} \\
5 & 52 & 0.008857 & \texttt{hollins.edu/admissions/info-request/info-request.cfm} \\
6 & 4023 & 0.008121 & \texttt{www1.hollins.edu/faculty/saloweyca/clas\%20395/Sculpture/sld001.htm} \\
7 & 43 & 0.007805 & \texttt{hollins.edu/admissions/apply/apply.htm} \\
8 & 27 & 0.006934 & \texttt{hollins.edu/admissions/admissions.htm} \\
9 & 3227 & 0.006924 & \texttt{www1.hollins.edu/faculty/saloweyca/clas\%20395/Sculpture/index.htm} \\
10 & 5254 & 0.006319 & \texttt{www1.hollins.edu/faculty/saloweyca/clas\%20395/Sculpture/tsld001.htm} \\
\bottomrule
\end{tabular}
\caption{The ten highest scores on Hollins ($m = 0.15$); the prefix \texttt{http://www.} is omitted}
\label{tab:hollins-top}
\end{table}




Every score is at least $m/n = 2.5\cdot10^{-5}$, as formula (3.2) predicts (the smallest score is $5.8\cdot10^{-5}$).

\exheading{11}{Ranking the Modified Web with the Matrix $\M$}{Ömer Özkan}

\exblock{Problem Statement}

We take the web of Figure~2.1 and add a page~5 that links to page~3, where page~3 also links to page~5.
We are asked to compute the new ranking by finding the eigenvector of $\M$ for $\lambda = 1$ with positive components summing to one, using $m = 0.15$.
This is the same web as in Exercise~1, but there the ranking was computed with $\A$; here we use $\M$.

\exblock{Web and Matrix Construction}

The new web has $n = 5$ pages and the links
\[
  1 \to 2,3,4 \qquad 2 \to 3,4 \qquad 3 \to 1,5 \qquad 4 \to 1,3 \qquad 5 \to 3 .
\]
\begin{itemize}
  \item Page~3 now has two outgoing links, so it splits its vote between page~1 and page~5 (entries $1/2$).
  \item Page~5 has only one link, so it gives its whole vote to page~3 (entry $1$).
  \item The other columns are the same as in Figure~2.1.
\end{itemize}
\[
  \A = \begin{pmatrix}
    0 & 0 & 1/2 & 1/2 & 0\\
    1/3 & 0 & 0 & 0 & 0\\
    1/3 & 1/2 & 0 & 1/2 & 1\\
    1/3 & 1/2 & 0 & 0 & 0\\
    0 & 0 & 1/2 & 0 & 0
  \end{pmatrix},
\]
\[
  \M = 0.85\,\A + 0.03\,\mathbf{1} =
  \begin{pmatrix}
    0.03 & 0.03 & 0.455 & 0.455 & 0.03\\
    0.3133 & 0.03 & 0.03 & 0.03 & 0.03\\
    0.3133 & 0.455 & 0.03 & 0.455 & 0.88\\
    0.3133 & 0.455 & 0.03 & 0.03 & 0.03\\
    0.03 & 0.03 & 0.455 & 0.03 & 0.03
  \end{pmatrix},
\]
where $\mathbf{1}$ is the $5\times5$ matrix of ones and $0.3133 = 0.85/3 + 0.03$.
Every column of $\M$ sums to one and every entry is at least $m/n = 0.03 > 0$.

\exblock{Why the Ranking Is Unique}

$\M$ is positive and column-stochastic, so by Lemma~3.2 of the paper $\dim\Vone{\M} = 1$.
There is exactly one eigenvector for $\lambda = 1$ with positive entries and sum one, and the power method converges to it (Proposition~5).

\exblock{Computation}

We run the power method of Section~\ref{sec:implementation} on formula (3.2) with $m = 0.15$, starting from the uniform vector, with $\texttt{tol} = 10^{-12}$.
It stops after 57 iterations with residual $6.6\cdot10^{-13}$, and gives
\[
  \vx \approx [0.237141,\ 0.097190,\ 0.348894,\ 0.138496,\ 0.178280] .
\]
\sumcheck{1.000}
The smallest component is $0.097 > 0$, as the theory says.
As a cross-check, the eigenvector of the dense matrix $\M$ computed with \texttt{numpy.linalg.eig} gives the same vector.

The ranking is $3 > 1 > 5 > 4 > 2$.
Table~\ref{tab:ex11} compares it with the ranking of the same web computed with $\A$ in Exercise~1.

\begin{table}[H]\centering
\begin{tabular}{cccc}
\toprule
\textbf{Page} & \textbf{Score with $\A$ (Ex.~1)} & \textbf{Score with $\M$, $m=0.15$} & \textbf{Rank} \\
\midrule
1 & 0.244898 & 0.237141 & 2 \\
2 & 0.081633 & 0.097190 & 5 \\
3 & 0.367347 & 0.348894 & 1 \\
4 & 0.122449 & 0.138496 & 4 \\
5 & 0.183673 & 0.178280 & 3 \\
\bottomrule
\end{tabular}
\caption{The modified web ranked with $\A$ and with $\M$}
\label{tab:ex11}
\end{table}

The order is the same in both cases.
With $\M$ the scores are a little closer to each other: the large scores go down and the small scores go up, because a part $m$ of every vote is shared equally among all pages.

\exblock{Conclusion}

The new ranking with $\M$ and $m = 0.15$ is
\[
  \vx \approx [0.2371,\ 0.0972,\ 0.3489,\ 0.1385,\ 0.1783], \qquad 3 > 1 > 5 > 4 > 2 .
\]
Page~3 is the most important page, as the owners of page~3 wanted (see Exercise~1).

\exheading{17}{The Choice of the Damping Factor $m$}{Ömer Özkan}

\exblock{Problem Statement}

We are asked how the value of $m$ should be chosen, and how this choice affects the rankings and the computation time.
We answer with an experiment (script \texttt{ex17.py}): we rank two webs for
$m \in \{0.05,\ 0.15,\ 0.25,\ 0.50,\ 0.85,\ 0.99\}$,
\begin{enumerate}
  \item the 4-page web of Figure~2.1,
  \item the Hollins dataset (6012 pages),
\end{enumerate}
and we record the number of iterations, the time and the ranking.
The power method starts from the uniform vector and stops when $\|\vx_{k+1}-\vx_k\|_1 < \texttt{tol} = 10^{-10}$.
We state the tolerance because the number of iterations has no meaning without it.

\exblock{Results: 4-Page Web (Figure 2.1)}

\subsubsection*{Convergence Speed}

Even on this tiny web the trend is clear: a larger $m$ needs fewer iterations (Table~\ref{tab:ex17-fig21}).

\begin{table}[H]\centering
\begin{tabular}{ccccccc}
\toprule
$\mathbf{m}$ & \textbf{Iterations} & \textbf{Time (s)} & $\mathbf{x_1}$ & $\mathbf{x_2}$ & $\mathbf{x_3}$ & $\mathbf{x_4}$ \\
\midrule
0.05 & 35 & 0.0006 & 0.380907 & 0.133120 & 0.289620 & 0.196353 \\
0.15 & 31 & 0.0004 & 0.368151 & 0.141809 & 0.287962 & 0.202078 \\
0.25 & 26 & 0.0003 & 0.354915 & 0.151229 & 0.285917 & 0.207940 \\
0.50 & 18 & 0.0002 & 0.320064 & 0.178344 & 0.278662 & 0.222930 \\
0.85 &  9 & 0.0002 & 0.269896 & 0.225995 & 0.261165 & 0.242944 \\
0.99 &  5 & 0.0001 & 0.251256 & 0.248338 & 0.250827 & 0.249579 \\
\bottomrule
\end{tabular}
\caption{Figure~2.1 for several values of $m$ ($\texttt{tol} = 10^{-10}$); the order is always $1>3>4>2$}
\label{tab:ex17-fig21}
\end{table}


\subsubsection*{Rankings}

The order $1 > 3 > 4 > 2$ never changes.
What changes is the distance between the scores: for $m = 0.05$ page~1 has $0.381$ and page~2 has $0.133$, for $m = 0.99$ all four scores are within $0.003$ of the uniform value $0.25$.
This is what formula (3.2) says: a part $m$ of every score is the same for all pages, so when $m$ grows the scores move towards $1/n$.

\exblock{Results: Hollins Dataset (6012 Pages)}

\subsubsection*{Convergence Speed}

On the large web the saving is important: going from $m = 0.15$ to $m = 0.85$ divides the number of iterations by almost forty, and the time goes from $2.9$~s to $0.15$~s (Table~\ref{tab:ex17-hollins}).

\begin{table}[h]\centering
\begin{tabular}{cccl}
\toprule
\textbf{m} & \textbf{Iterations} & \textbf{Time (s)} & \textbf{Top 5 pages} \\
\midrule
0.05 & 2850 & 36.33 & 4023, 3227, 4075, 5254, 3834 \\
0.15 &  407 &  2.91 & 2, 37, 38, 61, 52 \\
0.25 &  152 &  1.56 & 2, 37, 38, 61, 52 \\
0.50 &   39 &  0.45 & 2, 425, 37, 38, 52 \\
0.85 &   11 &  0.15 & 2, 425, 37, 38, 52 \\
0.99 &    5 &  0.07 & 2, 425, 37, 38, 52 \\
\bottomrule
\end{tabular}
\caption{Hollins for several values of $m$ (\texttt{tol} $= 10^{-10}$, uniform start)}
\label{tab:ex17-hollins}
\end{table}




\subsubsection*{Rankings}

The home page (page~2) is first for every $m \ge 0.15$, but the rest of the top five moves at $m = 0.50$.
Page~425 (a library page with a list of web links) is not in the top ten for $m = 0.15$ and 2nd for $m \ge 0.50$.
Looking at the data, page~425 has only 87 backlinks, but 57 of them come from pages with one or two links, so it receives their vote almost undivided.
For small $m$ these small pages have a low score and their votes are worth little; for large $m$ all scores are close to $1/n$, every voter is worth the same, and the undivided votes of page~425 push it up.
So on a real web a large $m$ does not only flatten the scores, it also changes the order.

For $m = 0.05$ the top five are pages of one course site (\texttt{faculty/saloweyca/clas 395}), which link only to each other: with dangling pages $\M$ loses vote at every step, and for very small $m$ the pages that keep their vote come first.

\exblock{Conclusion}

A larger $m$ makes the power method faster but moves every score towards $1/n$, so the ranking loses the information of the links.
On a small web the order does not change; on Hollins the top pages change from $m = 0.50$, and for a very small $m$ the ranking is decided by the dangling pages.
The value $m = 0.15$ used by Google is a reasonable compromise: the ranking follows the link structure, and the method converges in about four hundred iterations on a web of six thousand pages.

\clearpage
\part{Homework 3: Ten Exercises from the Paper}
\section*{About Part II}
\addcontentsline{toc}{section}{About Part II}

In this part we solve ten exercises of the paper: 1, 2, 3, 5, 6, 7, 8, 9, 12 and 13.
Each exercise is written as a short independent note with a problem statement, the solution and a conclusion.
The exercises with a computation (1, 2, 3, 12, 13) have a script in the repository with the same number; the others (5, 6, 7, 8, 9) are proofs.
Every numerical result was computed with the same solver as in Part~I.

Note on the numbering of the figures: Exercise~1 says ``the web of Figure~1'' and Exercise~3 says ``the web of Figure~2''; in the paper these figures are labelled Figure~2.1 and Figure~2.2.
We use the labels 2.1 and 2.2.

\exheading{1}{Can the Owners of Page 3 Boost Their Score?}{Ömer Özkan}

\exblock{Problem Statement}

In the web of Figure~2.1 the score of page~3, computed with formula (2.1), is lower than the score of page~1.
The owners of page~3 create a new page~5 that links to page~3, and page~3 also links to page~5.
We are asked whether this boosts the score of page~3 above the score of page~1.
Here the scores are computed with $\A$ only (formula (2.1)), so no damping.

\exblock{Before: the Web of Figure 2.1}

The paper gives the eigenvector $\vx = \frac{1}{31}(12,\,4,\,9,\,6)^T$:
\[
  x_1 \approx 0.387, \qquad x_3 \approx 0.290 .
\]
Page~1 is first and page~3 is third.

\exblock{After: Matrix Construction}

The new web has $n = 5$ pages.
\begin{itemize}
  \item Page~5 links only to page~3, so it gives its whole vote to page~3 (entry $1$ in column~5).
  \item Page~3 now links to page~1 and page~5, so it splits its vote (entries $1/2$ in column~3).
  \item Nothing else changes.
\end{itemize}
\[
  \A = \begin{pmatrix}
    0 & 0 & 1/2 & 1/2 & 0\\
    1/3 & 0 & 0 & 0 & 0\\
    1/3 & 1/2 & 0 & 1/2 & 1\\
    1/3 & 1/2 & 0 & 0 & 0\\
    0 & 0 & 1/2 & 0 & 0
  \end{pmatrix}.
\]

\exblock{Eigenvector Analysis}

We solve $\A\vx = \vx$.
The vector $\mathbf{v} = (12,\,4,\,18,\,6,\,9)^T$ satisfies it, as we check row by row:
\[
  \tfrac{18}{2}+\tfrac{6}{2} = 12,\quad
  \tfrac{12}{3} = 4,\quad
  \tfrac{12}{3}+\tfrac{4}{2}+\tfrac{6}{2}+9 = 18,\quad
  \tfrac{12}{3}+\tfrac{4}{2} = 6,\quad
  \tfrac{18}{2} = 9 .
\]
The web is strongly connected, so $\dim\Vone{\A} = 1$ and this is the only ranking (Exercise~10 of the paper; \texttt{numpy} confirms that $1$ is a simple eigenvalue).
Scaling to sum one gives
\[
  \vx = \tfrac{1}{49}(12,\,4,\,18,\,6,\,9)^T \approx [0.244898,\ 0.081633,\ 0.367347,\ 0.122449,\ 0.183673].
\]
\sumcheck{1.000}
Our power method with $m = 0$ (script \texttt{ex01.py}) gives the same numbers after 84 iterations.

Table~\ref{tab:ex1} compares the two webs. We also show the unscaled values, because they explain what happens.

\begin{table}[H]\centering
\begin{tabular}{ccccc}
\toprule
\textbf{Page} & \textbf{Before (unscaled)} & \textbf{Before} & \textbf{After (unscaled)} & \textbf{After} \\
\midrule
1 & 12 & 0.387097 & 12 & 0.244898 \\
2 & 4 & 0.129032 & 4 & 0.081633 \\
3 & 9 & 0.290323 & \textbf{18} & \textbf{0.367347} \\
4 & 6 & 0.193548 & 6 & 0.122449 \\
5 & --- & --- & 9 & 0.183673 \\
\midrule
sum & 31 & 1 & 49 & 1 \\
\bottomrule
\end{tabular}
\caption{Scores of Figure~2.1 before and after the new page~5}
\label{tab:ex1}
\end{table}

\begin{itemize}
  \item \textbf{Feedback loop.} Page~3 gives half of its score to page~5, and page~5 gives all of it back to page~3. So page~3 keeps what it had (the votes of pages 1, 2 and 4 are unchanged: $4 + 2 + 3 = 9$) and adds the vote of page~5, which is $18/2 = 9$. Its unscaled value doubles from 9 to 18.
  \item \textbf{Page 1 does not lose votes.} Its unscaled value stays 12. Before, it received the whole vote of page~3, that is $9$. Now it receives only half of the vote of page~3, but page~3 is twice as big, so this is again $18/2 = 9$; plus $3$ from page~4 as before. Its scaled score goes down only because the total grows from 31 to 49.
\end{itemize}

\exblock{Conclusion}

Yes. After the trick, page~3 has score $18/49 \approx 0.367$ and page~1 has $12/49 \approx 0.245$, so page~3 is now the most important page.
This shows that formula (2.1) can be manipulated by adding pages.

\exheading{2}{Dimension of $\Vone{\A}$ for a Web with Three Subwebs}{Zeynep Sude Ulaş}

\exblock{Problem Statement}

We are asked to construct a web consisting of three or more disconnected subwebs and to verify that $\dim\Vone{\A}$ equals (or exceeds) the number of subwebs.

\exblock{Web and Matrix Construction}

We take six pages divided into three subwebs,
\[
  W_1 = \{1,2\}, \qquad W_2 = \{3,4\}, \qquad W_3 = \{5,6\},
\]
where inside each subweb the two pages link to each other ($1 \leftrightarrow 2$, $3 \leftrightarrow 4$, $5 \leftrightarrow 6$) and there is no link between different subwebs.
Each page has exactly one outgoing link, so every non-zero entry of $\A$ is $1$:
\[
  \A = \begin{pmatrix}
    0&1&0&0&0&0\\ 1&0&0&0&0&0\\ 0&0&0&1&0&0\\ 0&0&1&0&0&0\\ 0&0&0&0&0&1\\ 0&0&0&0&1&0
  \end{pmatrix}
  = \begin{pmatrix} \A_1 & 0 & 0\\ 0 & \A_2 & 0\\ 0 & 0 & \A_3 \end{pmatrix},
  \qquad
  \A_1 = \A_2 = \A_3 = \begin{pmatrix} 0&1\\1&0 \end{pmatrix}.
\]
The matrix is block diagonal, one block for each subweb.

\exblock{Eigenvector Analysis}

Each block is column-stochastic and satisfies $\A_i (1,1)^T = (1,1)^T$, so $(1,1)^T$ is an eigenvector of each block for the eigenvalue $1$.
We extend these vectors to the full matrix by putting zeros in the other blocks:
\[
  \mathbf{w}_1 = (1,1,0,0,0,0)^T, \qquad
  \mathbf{w}_2 = (0,0,1,1,0,0)^T, \qquad
  \mathbf{w}_3 = (0,0,0,0,1,1)^T .
\]
Because $\A$ is block diagonal, $\A\mathbf{w}_i = \mathbf{w}_i$ for $i = 1,2,3$, so all three vectors are in $\Vone{\A}$.
They are linearly independent, because their non-zero entries are in different positions.
Therefore
\[
  \dim\Vone{\A} \ge 3 .
\]
In fact the dimension is exactly $3$: the eigenvalues of $\A$ computed by \texttt{numpy} (script \texttt{ex02.py}) are $1,-1,1,-1,1,-1$, so the eigenvalue $1$ appears three times.

\exblock{Conclusion}

For a web with three disconnected subwebs we found
\[
  \dim\Vone{\A} = 3 ,
\]
equal to the number of subwebs.
This confirms the statement of Section~2.2.1 of the paper: with $r$ subwebs, $\dim\Vone{\A} \ge r$, and so the ranking given by $\A$ is not unique.

\exheading{3}{A Connected Web Whose Ranking Is Still Not Unique}{Zeynep Sude Ulaş}

\exblock{Problem Statement}

We add a link from page~5 to page~1 in the web of Figure~2.2.
The new web, seen as an undirected graph, is connected.
We are asked for the dimension of $\Vone{\A}$.

\exblock{Web and Matrix Construction}

In Figure~2.2 the links are $1 \leftrightarrow 2$, $3 \leftrightarrow 4$ and $5 \to 3$, $5 \to 4$.
After adding $5 \to 1$, page~5 has three outgoing links, so it gives $1/3$ of its vote to each of the pages 1, 3 and 4:
\[
  \A = \begin{pmatrix}
    0 & 1 & 0 & 0 & 1/3\\
    1 & 0 & 0 & 0 & 0\\
    0 & 0 & 0 & 1 & 1/3\\
    0 & 0 & 1 & 0 & 1/3\\
    0 & 0 & 0 & 0 & 0
  \end{pmatrix}.
\]
Only column~5 changed with respect to the matrix of Figure~2.2.
Row~5 is still zero: nobody links to page~5.

\exblock{Eigenvector Analysis}

We solve $\A\vx = \vx$ by writing the five equations:
\[
  x_1 = x_2 + \tfrac{1}{3}x_5, \qquad
  x_2 = x_1, \qquad
  x_3 = x_4 + \tfrac{1}{3}x_5, \qquad
  x_4 = x_3, \qquad
  x_5 = 0 .
\]
The last equation gives $x_5 = 0$, so the new link carries no weight at all.
The system reduces to $x_1 = x_2$ and $x_3 = x_4$, and the general solution is
\[
  \vx = a\,(1,1,0,0,0)^T + b\,(0,0,1,1,0)^T, \qquad a,b \in \mathbb{R}.
\]
The two vectors are linearly independent, so $\dim\Vone{\A} = 2$.
The eigenvalues computed by \texttt{numpy} (script \texttt{ex03.py}) are $1,-1,1,-1,0$: the eigenvalue $1$ appears twice, which confirms the result.

\exblock{Conclusion}

\[
  \dim\Vone{\A} = 2 .
\]
Although the web is now connected as an undirected graph, the ranking given by $\A$ is still not unique.
The reason is the direction of the links: page~5 sends votes but receives none, so its score is zero and it cannot connect the two subwebs $\{1,2\}$ and $\{3,4\}$ in the equations.
Being connected as an undirected graph is therefore not enough for $\dim\Vone{\A} = 1$; the paper's condition of a \emph{strongly} connected web (Exercise~10) is what matters.

\exheading{5}{Importance Score of a Page with No Backlinks}{Zeynep Sude Ulaş}

\exblock{Problem Statement}

We are asked to prove that, in any web, the importance score of a page with no backlinks is zero.
Here the score is defined by formula (2.1), that is, by the matrix $\A$.

\exblock{Proof}

By formula (2.1) the score of page $k$ is
\[
  x_k = \sum_{j \in L_k} \frac{x_j}{n_j},
\]
where $L_k$ is the set of pages that link to page $k$.
If page $k$ has no backlinks, then $L_k = \varnothing$.
A sum over the empty set is zero, so
\[
  x_k = 0 .
\]
In matrix language: page $k$ has no backlinks means that row $k$ of $\A$ is zero, so $x_k = (\A\vx)_k = 0$ for every $\vx \in \Vone{\A}$.

\exblock{Conclusion}

With formula (2.1), a page that nobody links to has importance score $x_k = 0$, whatever the rest of the web looks like.
Exercise~9 shows that with the matrix $\M$ this score becomes $m/n > 0$.

\exheading{6}{The Scores Do Not Depend on the Numbering of the Pages}{Zeynep Sude Ulaş}

\exblock{Problem Statement}

We are asked to prove that the way the pages are numbered has no effect on the importance scores.
Let $\A$ be the link matrix of a web with $n$ pages, and let $\tilde{\A}$ be the link matrix of the same web after the indices of pages $i$ and $j$ are exchanged.
Let $\mathbf{P}$ be the permutation matrix obtained by exchanging rows $i$ and $j$ of the identity matrix.

\exblock{Proof}

\subsubsection*{Step 1: $\tilde{\A} = \mathbf{P}\A\mathbf{P}$}

Multiplying on the left, $\mathbf{P}\A$, exchanges rows $i$ and $j$ of $\A$.
Multiplying on the right, $\A\mathbf{P}$, exchanges columns $i$ and $j$.
When we rename page $i$ as page $j$ and vice versa, the row of page $i$ becomes the row of page $j$ and the column of page $i$ becomes the column of page $j$.
So both the rows and the columns are exchanged, which is exactly $\mathbf{P}\A\mathbf{P}$.
Also $\mathbf{P}^2 = \mathbf{I}$, because exchanging twice gives the identity, and therefore $\mathbf{P}^{-1} = \mathbf{P}$.

\subsubsection*{Step 2: $\mathbf{y} = \mathbf{P}\vx$ is an eigenvector of $\tilde{\A}$}

Suppose $\A\vx = \lambda\vx$ and define $\mathbf{y} = \mathbf{P}\vx$. Then
\[
  \tilde{\A}\mathbf{y} = (\mathbf{P}\A\mathbf{P})(\mathbf{P}\vx) = \mathbf{P}\A(\mathbf{P}^2)\vx = \mathbf{P}\A\vx = \mathbf{P}(\lambda\vx) = \lambda\,\mathbf{P}\vx = \lambda\mathbf{y}.
\]
So $\mathbf{y}$ is an eigenvector of $\tilde{\A}$ with the same eigenvalue $\lambda$.

\subsubsection*{Step 3: the scores are unchanged}

Take $\lambda = 1$ and let $\vx$ be the score vector, $\A\vx = \vx$.
Then $\mathbf{y} = \mathbf{P}\vx$ is the score vector of the renamed web.
But $\mathbf{P}\vx$ is just $\vx$ with the entries $i$ and $j$ exchanged: the numbers are the same, only their positions move.
So page $i$ (now called page $j$) has the same score as before.
Any permutation of the $n$ indices is a product of such exchanges, so the same holds for any renumbering: if $\mathbf{Q}$ is a permutation matrix, $\hat{\A} = \mathbf{Q}\A\mathbf{Q}^{-1}$ and $\hat{\A}(\mathbf{Q}\vx) = \mathbf{Q}\A\vx = \mathbf{Q}\vx$.

\exblock{Conclusion}

Renumbering the pages permutes the entries of the score vector in the same way, and the scores themselves do not change.
The importance scores are therefore independent of the numbering of the pages.

\exheading{7}{The Matrix $\M$ Is Column-Stochastic}{Zeynep Sude Ulaş}

\exblock{Problem Statement}

We are asked to prove that if $\A$ is an $n \times n$ column-stochastic matrix and $0 \le m \le 1$, then
\[
  \M = (1-m)\A + m\Smat
\]
is also column-stochastic. Here $\Smat$ is the matrix with all entries equal to $1/n$.

\exblock{Proof}

A matrix is column-stochastic when all its entries are non-negative and every column sums to one.
Both $\A$ and $\Smat$ have these properties:
\[
  a_{ij} \ge 0, \quad s_{ij} = \tfrac{1}{n} \ge 0, \qquad
  \sum_{i=1}^n a_{ij} = 1, \quad \sum_{i=1}^n s_{ij} = n \cdot \tfrac{1}{n} = 1 \quad \text{for every } j.
\]

\textbf{Non-negative entries.} The entries of $\M$ are $M_{ij} = (1-m)a_{ij} + m s_{ij}$.
Since $0 \le m \le 1$, both $1-m \ge 0$ and $m \ge 0$, so $M_{ij}$ is a sum of two non-negative numbers and $M_{ij} \ge 0$.

\textbf{Column sums.} For every column $j$,
\[
  \sum_{i=1}^n M_{ij} = (1-m)\sum_{i=1}^n a_{ij} + m\sum_{i=1}^n s_{ij} = (1-m)\cdot 1 + m \cdot 1 = 1 .
\]

\exblock{Conclusion}

All entries of $\M$ are non-negative and every column of $\M$ sums to one, so $\M$ is column-stochastic.
By Proposition~1 of the paper, $1$ is then an eigenvalue of $\M$.
If moreover $m > 0$, every entry satisfies $M_{ij} \ge m/n > 0$, so $\M$ is positive and Lemma~3.2 gives $\dim\Vone{\M} = 1$.

\exheading{8}{The Product of Two Column-Stochastic Matrices}{Ömer Özkan}

\exblock{Problem Statement}

We are asked to show that the product of two column-stochastic matrices is also column-stochastic.

\exblock{Definitions}

Let $\mathbf{B}$ and $\mathbf{C}$ be $n \times n$ column-stochastic matrices:
\[
  b_{ik} \ge 0,\quad c_{kj} \ge 0, \qquad \sum_{i=1}^n b_{ik} = 1,\quad \sum_{k=1}^n c_{kj} = 1 \qquad \text{for all } i,j,k .
\]
The entries of the product are $(\mathbf{BC})_{ij} = \sum_{k=1}^n b_{ik}c_{kj}$.

\exblock{Proof}

\textbf{Non-negative entries.} Each term $b_{ik}c_{kj}$ is a product of two non-negative numbers, so $(\mathbf{BC})_{ij} \ge 0$.

\textbf{Column sums.} Fix a column $j$ and sum over $i$. We exchange the order of the two sums:
\[
  \sum_{i=1}^n (\mathbf{BC})_{ij}
  = \sum_{i=1}^n \sum_{k=1}^n b_{ik}c_{kj}
  = \sum_{k=1}^n c_{kj} \left( \sum_{i=1}^n b_{ik} \right)
  = \sum_{k=1}^n c_{kj} \cdot 1
  = 1 .
\]
In the third step we used that every column $k$ of $\mathbf{B}$ sums to one, and in the last step that column $j$ of $\mathbf{C}$ sums to one.

\exblock{Why This Matters}

By induction, every power $\A^k$ of a column-stochastic matrix is column-stochastic, and so is $\M^k$.
This is used in the last part of Exercise~10 of the paper, where the matrix $\mathbf{B} = \frac{1}{n}(\mathbf{I} + \A + \dots + \A^{n-1})$ must be column-stochastic.
It also means that in the power method the vector $\vx_k = \M^k\vx_0$ keeps the sum $\sum_i (\vx_k)_i = 1$ at every step, so no rescaling is needed in exact arithmetic.

\exblock{Conclusion}

If $\mathbf{B}$ and $\mathbf{C}$ are column-stochastic, then $\mathbf{BC}$ has non-negative entries and columns that sum to one: $\mathbf{BC}$ is column-stochastic.

\exheading{9}{A Page with No Backlinks Gets the Score $m/n$}{Ömer Özkan}

\exblock{Problem Statement}

We are asked to show that, with formula (3.2), a page with no backlinks gets the importance score $m/n$.
This is the version with damping of Exercise~5, where the same page got the score $0$.

\exblock{Definitions}

Formula (3.2) is
\[
  \vx = (1-m)\A\vx + m\vs ,
\]
where $\vs$ is the vector with all entries $1/n$ and $\vx$ is scaled so that $\sum_i x_i = 1$.
Page $i$ has no backlinks when no page links to it, that is, when row $i$ of $\A$ is zero.

\exblock{Proof}

Take row $i$ of formula (3.2):
\[
  x_i = (1-m)\sum_{j=1}^n A_{ij}x_j + \frac{m}{n} .
\]
If page $i$ has no backlinks, then $A_{ij} = 0$ for every $j$, so the sum is zero and
\[
  x_i = \frac{m}{n} .
\]
The result does not depend on the rest of the web, only on $m$ and $n$.
Also, since $\A\vx \ge 0$, every page gets at least $m/n$: this is the smallest possible score with the matrix $\M$.

\exblock{Check with Our Computations}

\begin{itemize}
  \item In Exercise~12, page~6 has no backlinks, $n = 6$, $m = 0.15$, and we compute $x_6 = 0.025000 = 0.15/6$.
  \item In Exercise~13, page~5 has no backlinks, $n = 5$, and we compute $x_5 = 0.030000 = 0.15/5$.
  \item In Example~2 of the paper (Figure~2.2), page~5 has no backlinks and the paper gives $x_5 = 0.03 = 0.15/5$.
  \item On Hollins every score is at least $m/n = 2.5\cdot10^{-5}$ (Section~\ref{sec:hollins}).
\end{itemize}

\exblock{A Remark on the Paper}

In Section~3.1 the paper writes ``a web page with no backlinks (a dangling node)''.
These are two different things: a page with no backlinks has a zero \emph{row} in $\A$ (nobody links to it), while a dangling node has a zero \emph{column} (it links to nobody).
Exercise~9 is about the zero row.
For example page~5 of Figure~2.2 has no backlinks but is not dangling, since it links to pages 3 and 4.

\exblock{Conclusion}

With formula (3.2) a page with no backlinks has importance score exactly
\[
  x_i = \frac{m}{n} ,
\]
instead of $0$ as with formula (2.1). The damping term gives every page a small base score.

\exheading{12}{A Sixth Page That Links to Everyone}{Ömer Özkan}

\exblock{Problem Statement}

We add to the web of Exercise~11 a sixth page that links to every other page, but to which no page links.
We are asked to rank the pages with $\A$ and then with $\M$ ($m = 0.15$), and to compare the results.

\exblock{Web and Matrix Construction}

The new web has $n = 6$ pages.
\begin{itemize}
  \item Page~6 has five outgoing links, so column~6 has the entry $1/5$ in rows 1 to 5.
  \item Nobody links to page~6, so row~6 is zero.
  \item Page~6 is \emph{not} a dangling node: it has outgoing links, so its column sums to one. What it lacks is backlinks.
\end{itemize}
\[
  \A = \begin{pmatrix}
    0 & 0 & 1/2 & 1/2 & 0 & 1/5\\
    1/3 & 0 & 0 & 0 & 0 & 1/5\\
    1/3 & 1/2 & 0 & 1/2 & 1 & 1/5\\
    1/3 & 1/2 & 0 & 0 & 0 & 1/5\\
    0 & 0 & 1/2 & 0 & 0 & 1/5\\
    0 & 0 & 0 & 0 & 0 & 0
  \end{pmatrix}.
\]

\exblock{Ranking with $\A$ and with $\M$}

Both rankings are computed with the same solver (script \texttt{ex12.py}): $m = 0$ gives formula (2.1), $m = 0.15$ gives formula (3.2).
Table~\ref{tab:ex12} shows that the order of the pages is the same in both cases, and that only page~6 changes its nature: from $0$ to $m/n$.

\begin{table}[H]\centering
\begin{tabular}{cccc}
\toprule
\textbf{Page} & \textbf{Score with $\A$} & \textbf{Score with $\M$, $m=0.15$} & \textbf{Rank} \\
\midrule
1 & 0.244898 & 0.231212 & 2 \\
2 & 0.081633 & 0.094760 & 5 \\
3 & 0.367347 & 0.340172 & 1 \\
4 & 0.122449 & 0.135033 & 4 \\
5 & 0.183673 & 0.173823 & 3 \\
6 & \textbf{0.000000} & \textbf{0.025000} & 6 \\
\bottomrule
\end{tabular}
\caption{The six-page web ranked with $\A$ and with $\M$}
\label{tab:ex12}
\end{table}

\begin{itemize}
  \item \textbf{With $\A$.} Page~6 has no backlinks, so $x_6 = 0$ (Exercise~5). Its column then multiplies $x_6 = 0$ and has no effect: the other five scores are exactly those of Exercise~1, $\frac{1}{49}(12,4,18,6,9)$.
  \item \textbf{With $\M$.} Page~6 gets exactly $x_6 = m/n = 0.15/6 = 0.025$ (Exercise~9). The other five scores are $0.975$ times the scores of Exercise~11, for example $0.975 \times 0.237141 = 0.231212$.
  The reason is simple: page~6 gives $\frac{1}{5}x_6$ to every page, the same amount to all, just like the term $m\vs$. So pages 1 to 5 satisfy the same equations as in Exercise~11 with a different constant, and the solution of such a system is proportional to the constant. The proportions do not change, and the sum of the five scores must be $1 - 0.025 = 0.975$.
\end{itemize}

\exblock{Conclusion}

The order is $3 > 1 > 5 > 4 > 2 > 6$ with both matrices.
With $\A$ the new page scores $0$; with $\M$ it scores $m/n = 0.025$, and the other pages keep their relative importance.
Linking to everyone does not help a page that nobody links to: the score comes from the backlinks, not from the outgoing links.

\exheading{13}{Ranking a Web with Two Subwebs Using $\M$}{Zeynep Sude Ulaş}

\exblock{Problem Statement}

We are asked to construct a web consisting of two or more subwebs and to determine the ranking given by formula (3.1), that is, by the matrix $\M = (1-m)\A + m\Smat$.
We use $m = 0.15$.

\exblock{Web and Matrix Construction}

We take five pages in two disconnected subwebs,
\[
  W_1 = \{1,2\}, \qquad W_2 = \{3,4,5\},
\]
with the links $1 \to 2$, $2 \to 1$, $3 \to 4$, $4 \to 3$ and $5 \to 3$.
There is no link between $W_1$ and $W_2$, and every page has exactly one outgoing link, so
\[
  \A = \begin{pmatrix}
    0&1&0&0&0\\ 1&0&0&0&0\\ 0&0&0&1&1\\ 0&0&1&0&0\\ 0&0&0&0&0
  \end{pmatrix},
  \qquad
  \M = 0.85\,\A + 0.03\,\mathbf{1} = \begin{pmatrix}
    0.03&0.88&0.03&0.03&0.03\\
    0.88&0.03&0.03&0.03&0.03\\
    0.03&0.03&0.03&0.88&0.88\\
    0.03&0.03&0.88&0.03&0.03\\
    0.03&0.03&0.03&0.03&0.03
  \end{pmatrix},
\]
where $\mathbf{1}$ is the $5 \times 5$ matrix of ones and $0.03 = m/n$.

\exblock{Why $\A$ Is Not Enough}

The web has two subwebs, so $\dim\Vone{\A} = 2$ (the eigenvalues of $\A$ are $1,-1,1,-1,0$, script \texttt{ex13.py}).
Any combination $a\,(1,1,0,0,0)^T + b\,(0,0,1,1,0)^T$ is an eigenvector, so $\A$ gives no unique ranking and cannot compare a page of $W_1$ with a page of $W_2$.

\exblock{Ranking with $\M$}

$\M$ is positive, so $\dim\Vone{\M} = 1$ (Lemma~3.2) and the ranking is unique.
Solving $\M\vx = \vx$ with $\sum_i x_i = 1$ gives
\[
  \vx = \left[\tfrac{1}{5},\ \tfrac{1}{5},\ \tfrac{54}{185},\ \tfrac{1029}{3700},\ \tfrac{3}{100}\right]
  \approx [0.200000,\ 0.200000,\ 0.291892,\ 0.278108,\ 0.030000].
\]
\sumcheck{1.000}
Our power method gives the same vector after 165 iterations ($\texttt{tol} = 10^{-12}$).

\begin{itemize}
  \item Pages 1 and 2 tie exactly: the subweb $1 \leftrightarrow 2$ is symmetric, and formula (3.2) gives $x_1 = 0.85\,x_2 + 0.03$ and $x_2 = 0.85\,x_1 + 0.03$, so $x_1 = x_2 = 0.2$.
  \item Page 5 has no backlinks, so its score is exactly $m/n = 0.03$ (Exercise~9).
  \item Page 3 is above page 4 because it receives the votes of both page 4 and page 5, while page 4 receives only the vote of page 3.
\end{itemize}

\exblock{Conclusion}

The ranking given by formula (3.1) is
\[
  3 > 4 > 1 = 2 > 5 .
\]
Although the web consists of two disconnected subwebs, the matrix $\M$ produces a unique ranking and makes the pages of different subwebs comparable.


\clearpage
\begin{thebibliography}{9}
\bibitem{bryanleise} K.~Bryan and T.~Leise,
  \emph{The \$25,000,000,000 Eigenvector: The Linear Algebra Behind Google},
  SIAM Review, 48 (2006), pp.~569--581.
\bibitem{repo} Ö.~Özkan and Z.~S.~Ulaş, \emph{pagerank\_project} (source code),
  \url{https://github.com/omerzkan/pagerank_project}.
\end{thebibliography}

\end{document}
